\phantomsection
\part{論理}

\begin{abstract}
この章では，数学の議論を支える「論理」を整理する．まず命題論理の記号と真理値表を確認し，交換律・分配律・ド・モルガンの法則などの基本法則を述べる．含意の真偽をめぐる注意と，逆・裏・対偶の関係もここで確認する．次に，変数を含む文を扱う述語論理を導入し，全称命題・存在命題および量化子の入れ子について注意しよう．
\end{abstract}

\phantomsection
\section{命題論理\index{めいだいろんり@命題論理}}

命題論理では，命題を結合するために記号を用いる．代表的な論理記号を列挙してみよう：

\begin{table}[ht]
  \centering
  \caption{命題論理の記号}
  \begin{tabular}{c|c}
    \hline
    記号                & 意味                       \\
    \hline
    $\land$           & かつ（論理積\index{ろんりせき@論理積}） \\
    $\lor$            & または（論理和\index{ろんりわ@論理和}） \\
    $\lnot$           & 否定\index{ひてい@否定}         \\
    $\to$             & ならば（含意\index{がんい@含意}）    \\
    $\leftrightarrow$ & 同値 \index{どうち@同値} \\
    \hline
  \end{tabular}
\end{table}

これらの論理記号を用いて，命題を結合し，複雑な論理式を構成することができる．
以下に，基本的な命題の真理値表を載せよう．

\begin{table}[htbp]
\centering
\small

%================================================
% 上段：3段の表
%================================================
\begin{minipage}[t]{0.45\linewidth}
\centering
\caption{否定命題の真理値表}
\label{tab:否定命題の真理値表}
\begin{tabular}{|c|c|} \hline
  $P$          & $\lnot P$    \\ \hline
  $\mathrm{T}$ & $\mathrm{F}$ \\ \hline
  $\mathrm{F}$ & $\mathrm{T}$ \\ \hline
\end{tabular}
\end{minipage}
\hfill
\begin{minipage}[t]{0.45\linewidth}
\centering
\caption{恒真命題と矛盾命題の真理値表}
\label{tab:恒真命題と矛盾命題の真理値表}
\begin{tabular}{|c|c|c|c|} \hline
  $P$          & $\lnot P$    & $P \lor \lnot P$ & $P \land \lnot P$ \\ \hline
  $\mathrm{T}$ & $\mathrm{F}$ & $\mathrm{T}$     & $\mathrm{F}$      \\ \hline
  $\mathrm{F}$ & $\mathrm{T}$ & $\mathrm{T}$     & $\mathrm{F}$      \\ \hline
\end{tabular}
\end{minipage}

\vspace{6mm}

%================================================
% 下段：5段の表
%================================================
\begin{minipage}[t]{0.32\linewidth}
\centering
\caption{論理和（積）の真理値表}
\label{tab:論理和（積）の真理値表}
\begin{tabular}{|c|c|c|c|} \hline
  $P$          & $Q$          & $P \land Q$  & $P \lor Q$   \\ \hline
  $\mathrm{T}$ & $\mathrm{T}$ & $\mathrm{T}$ & $\mathrm{T}$ \\ \hline
  $\mathrm{T}$ & $\mathrm{F}$ & $\mathrm{F}$ & $\mathrm{T}$ \\ \hline
  $\mathrm{F}$ & $\mathrm{T}$ & $\mathrm{F}$ & $\mathrm{T}$ \\ \hline
  $\mathrm{F}$ & $\mathrm{F}$ & $\mathrm{F}$ & $\mathrm{F}$ \\ \hline
\end{tabular}
\end{minipage}
\hfill
\begin{minipage}[t]{0.32\linewidth}
\centering
\caption{「含意」の真理値表}
\label{tab:含意の真理値表}
\begin{tabular}{|c|c|c|} \hline
  $P$          & $Q$          & $P \to Q$    \\ \hline
  $\mathrm{T}$ & $\mathrm{T}$ & $\mathrm{T}$ \\ \hline
  $\mathrm{T}$ & $\mathrm{F}$ & $\mathrm{F}$ \\ \hline
  $\mathrm{F}$ & $\mathrm{T}$ & $\mathrm{T}$ \\ \hline
  $\mathrm{F}$ & $\mathrm{F}$ & $\mathrm{T}$ \\ \hline
\end{tabular}
\end{minipage}
\hfill
\begin{minipage}[t]{0.32\linewidth}
\centering
\caption{「同値」の真理値表}
\label{tab:同値の真理値表}
\begin{tabular}{|c|c|c|} \hline
  $P$          & $Q$          & $P \leftrightarrow Q$ \\ \hline
  $\mathrm{T}$ & $\mathrm{T}$ & $\mathrm{T}$          \\ \hline
  $\mathrm{T}$ & $\mathrm{F}$ & $\mathrm{F}$          \\ \hline
  $\mathrm{F}$ & $\mathrm{T}$ & $\mathrm{F}$          \\ \hline
  $\mathrm{F}$ & $\mathrm{F}$ & $\mathrm{T}$          \\ \hline
\end{tabular}
\end{minipage}

\end{table}

上記の真理値表によって，命題論理の記号を定義する．
以下，左辺$P$と右辺の$Q$の真理値が等しいことを$P \equiv Q$と表し，この記号をもって「同値」を表すことにする．

\begin{prop}{}{}
  命題$P$，$Q$，$R$に対して，次の式が成り立つ：
  \begin{description}
    \item[交換律]
          \[
            P \land Q \equiv Q \land P,\quad
            P \lor Q \equiv Q \lor P
          \]
    \item[結合律]
          \[
            (P \land Q) \land R \equiv P \land (Q \land R),\quad
            (P \lor Q) \lor R \equiv P \lor (Q \lor R)
          \]
    \item[分配律]
          \[
            P \land (Q \lor R) \equiv (P \land Q) \lor (P \land R),\quad
            P \lor (Q \land R) \equiv (P \lor Q) \land (P \lor R)
          \]
    \item[ド・モルガンの法則]
          \[
            \lnot (P \land Q) \equiv \lnot P \lor \lnot Q,\quad
            \lnot (P \lor Q) \equiv \lnot P \land \lnot Q
          \]
  \end{description}
\end{prop}

\begin{prop}{}{}
  命題$P$，$Q$に対して，次の式が成り立つ：
  \begin{description}
    \item[含意の等価表現] 
          \[
            P \to Q \equiv \lnot P \lor Q
          \]
    \item[同値の等価表現]
          \[
            P \leftrightarrow Q \equiv (P \to Q) \land (Q \to P)
          \]
  \end{description}
\end{prop}

\phantomsection
\subsection{含意の真偽と逆・裏・対偶}

含意$P \to Q$は，慣れないうちは真偽の判定にとまどうことが多い．そこで，表\ref{tab:含意の真理値表}にしたがって，$P$と$Q$の真偽の組合せ$4$通りをすべて列挙して確認しておこう．

\begin{center}
  \begin{tabular}{|c|c|c|l|} \hline
    $P$ & $Q$ & $P \to Q$   & 備考                                   \\ \hline
    真  & 真  & 真          & 約束は守られている                     \\ \hline
    真  & 偽  & \textbf{偽} & 偽となるのは，この場合\textbf{だけ}    \\ \hline
    偽  & 真  & 真          & $P$が成り立たないので，約束は破られない \\ \hline
    偽  & 偽  & 真          & 同上                                   \\ \hline
  \end{tabular}
\end{center}

とくに注意したいのが3行目，すなわち「$P$でない，かつ$Q$である」（$\lnot P \land Q$）という場合である．このときにも$P \to Q$は真とされる．
これは，含意$P \to Q$が「$P$が成り立つ場合に，$Q$も成り立つ」ということのみを主張しており，$P$が成り立たない場合については\textbf{なにも主張していない}からである．
「約束が破られるのは，$P$が成り立ったのに$Q$が成り立たなかったときだけ」と考えるとよい\footnote{このように，仮定$P$が偽であるために$P \to Q$が無条件に真となることを，\emph{空虚な真}（vacuous truth）ということがある．}．

\phantomsection
\begin{example}{含意の真偽}{含意の真偽}
  実数$x$に関する命題「$x > 2$ならば$x > 0$」は真である．このことを，いくつかの$x$の値で確認してみよう．
  \begin{center}
    \begin{tabular}{|c|c|c|l|} \hline
      $x$  & $x > 2$（$P$） & $x > 0$（$Q$） & 約束は守られているか                     \\ \hline
      $3$  & 真             & 真             & 守られている                             \\ \hline
      $1$  & 偽             & 真             & 仮定が成り立たないので，破られてはいない \\ \hline
      $-1$ & 偽             & 偽             & 同上                                     \\ \hline
    \end{tabular}
  \end{center}
  $x = 1$の場合が「$P$でない，かつ$Q$である」場合にあたるが，$x > 2$が成り立っていない以上，約束が破られたわけではない．
  どの$x$についても「$x>2$が成り立つのに$x>0$が成り立たない」ということは起こらないので，この命題は真である．
\end{example}

次に，含意$P \to Q$から作られる命題に名前をつけよう．

\phantomsection
\begin{definition}{否定・逆・裏・対偶}{否定・逆・裏・対偶}
  命題$P \to Q$に対して，以下を定義する．
  \begin{center}
    \begin{tabular}{|l|c|} \hline
      名称                      & 命題                  \\ \hline
      否定\index{ひてい@否定}   & $\lnot ( P \to Q )$   \\ \hline
      逆\index{ぎゃく@逆}       & $Q \to P$             \\ \hline
      裏\index{うら@裏}         & $\lnot P \to \lnot Q$ \\ \hline
      対偶\index{たいぐう@対偶} & $\lnot Q \to \lnot P$ \\ \hline
    \end{tabular}
  \end{center}
\end{definition}

まず，含意の否定について次が成り立つ．

\begin{prop}{含意の否定}{含意の否定}
  命題$P$，$Q$に対して，
  \[
    \lnot ( P \to Q ) \equiv P \land \lnot Q
  \]
  である．
\end{prop}

\begin{tleftbar}
  \begin{proof}
    含意の等価表現$P \to Q \equiv \lnot P \lor Q$とド・モルガンの法則を用いると，
    \[
      \lnot ( P \to Q )
      \equiv \lnot ( \lnot P \lor Q )
      \equiv \lnot \lnot P \land \lnot Q
      \equiv P \land \lnot Q
    \]
    である．これが証明すべきことであった．
  \end{proof}
\end{tleftbar}

すなわち，「$P$ならば$Q$」の否定は「$P$であるのに$Q$でない」である．「$P$ならば$Q$でない」（$P \to \lnot Q$）ではないことに注意しよう．

続いて，逆・裏・対偶と元の命題の関係を，真理値表を用いて調べよう．

\begin{prop}{対偶との同値}{対偶との同値}
  命題$P$，$Q$に対して，
  \[
    P \to Q \equiv \lnot Q \to \lnot P
  \]
  である．すなわち，命題とその対偶の真偽は一致する．同様に，逆$Q \to P$と裏$\lnot P \to \lnot Q$の真偽も一致する．

  一方，$P \to Q$と逆$Q \to P$の真偽は，一般には一致\textbf{しない}．
\end{prop}

\begin{tleftbar}
  \begin{proof}
    $P$，$Q$の真偽の組合せ$4$通りすべてについて，各命題の真偽を列挙する：
    \begin{center}
      \begin{tabular}{|c|c||c|c||c|c|c|c|} \hline
        $P$          & $Q$          & $\lnot P$    & $\lnot Q$    & $P \to Q$    & $\lnot Q \to \lnot P$ & $Q \to P$    & $\lnot P \to \lnot Q$ \\ \hline
        $\mathrm{T}$ & $\mathrm{T}$ & $\mathrm{F}$ & $\mathrm{F}$ & $\mathrm{T}$ & $\mathrm{T}$          & $\mathrm{T}$ & $\mathrm{T}$          \\ \hline
        $\mathrm{T}$ & $\mathrm{F}$ & $\mathrm{F}$ & $\mathrm{T}$ & $\mathrm{F}$ & $\mathrm{F}$          & $\mathrm{T}$ & $\mathrm{T}$          \\ \hline
        $\mathrm{F}$ & $\mathrm{T}$ & $\mathrm{T}$ & $\mathrm{F}$ & $\mathrm{T}$ & $\mathrm{T}$          & $\mathrm{F}$ & $\mathrm{F}$          \\ \hline
        $\mathrm{F}$ & $\mathrm{F}$ & $\mathrm{T}$ & $\mathrm{T}$ & $\mathrm{T}$ & $\mathrm{T}$          & $\mathrm{T}$ & $\mathrm{T}$          \\ \hline
      \end{tabular}
    \end{center}
    表より，$P \to Q$の列と$\lnot Q \to \lnot P$の列は$4$通りすべての場合で一致しているので，$P \to Q \equiv \lnot Q \to \lnot P$である．同様に，$Q \to P$の列と$\lnot P \to \lnot Q$の列も一致している．

    一方，$P \to Q$の列と$Q \to P$の列は，$2$行目・$3$行目で真偽が異なっている．よって$P \to Q$と$Q \to P$は同値でない．これが証明すべきことであった．
  \end{proof}
\end{tleftbar}

\begin{remark}[対偶と背理法]
  命題とその対偶の真偽が一致することから，$P \to Q$を直接示す代わりに，対偶$\lnot Q \to \lnot P$を示してもよい．これを\emph{対偶法}\index{たいぐうほう@対偶法}（対偶による証明）という．のちに紹介する背理法と並んで，重要な証明法のひとつである．
\end{remark}

\phantomsection
\begin{example}{逆・裏・対偶の例}{逆・裏・対偶の例}
  実数$x$に関する命題「$x > 2$ならば$x > 0$」（真）について，
  \begin{center}
    \begin{tabular}{|l|c|l|} \hline
      名称 & 命題                                  & 真偽                 \\ \hline
      逆   & 「$x > 0$ならば$x > 2$」              & 偽（反例：$x = 1$）  \\ \hline
      裏   & 「$x \leqq 2$ならば$x \leqq 0$」      & 偽（反例：$x = 1$）  \\ \hline
      対偶 & 「$x \leqq 0$ならば$x \leqq 2$」      & 真（元の命題と同じ） \\ \hline
    \end{tabular}
  \end{center}
  このように，元の命題が真であっても，その逆・裏は偽でありうる．
\end{example}

\begin{remark}[「逆」の誤用に注意]
  「逆」という言葉は，日常語では「向きを反対にすること」全般を指すため，数学の議論において誤用されやすい．たとえば，三角関数の加法定理
  \[
    \sin (\alpha + \beta) = \sin \alpha \cos \beta + \cos \alpha \sin \beta
  \]
  を右辺から左辺へ向けて用いて，$\sin \alpha \cos \beta + \cos \alpha \sin \beta$を$\sin (\alpha + \beta)$にまとめる変形を「加法定理の逆」とよぶのは誤りである．等式$A = B$と$B = A$は同じ内容を表しており，等式をどちら向きに用いても「加法定理そのもの」だからである．\deref{否定・逆・裏・対偶}で定義したとおり，「逆」とは含意$P \to Q$の仮定と結論を入れ替えた$Q \to P$のことをいうのであって，等式を用いる向きを指す言葉ではない．

  一方，加法定理を用いて\textbf{得られる含意}については，通常の意味で逆を考えることができる．たとえば，実数$\alpha$，$\beta$に関する命題
  \[
    \alpha + \beta = \pi
    \quad \Longrightarrow \quad
    \sin \alpha \cos \beta + \cos \alpha \sin \beta = 0
  \]
  は加法定理から従う．実際，仮定$\alpha + \beta = \pi$のもとで
  \[
    \sin \alpha \cos \beta + \cos \alpha \sin \beta
    = \sin (\alpha + \beta)
    = \sin \pi
    = 0
  \]
  となる．この命題の逆
  \[
    \sin \alpha \cos \beta + \cos \alpha \sin \beta = 0
    \quad \Longrightarrow \quad
    \alpha + \beta = \pi
  \]
  は\textbf{偽}である．実際，$\alpha = \beta = 0$とすれば，仮定の等式は成り立つが，$\alpha + \beta = 0 \ne \pi$である．
\end{remark}


\phantomsection
\section{述語論理\index{じゅつごろんり@述語論理}}

\exref{命題でない文}の\ref{enu:命題でない文4}において，「$x^2 >4$」という文について考えた．この文は命題ではないが，たとえば$x=3$を代入したときは真，$x=0$を代入したときは偽である．
つまり，$x$になにかしら値を代入すると，この文は真偽が決まり，命題となるのである．

このような「値が決まっていない変数を含み，変数になにかを代入すると真偽が判定できる」という性質を持つ文を\emph{述語}\index{じゅつご@述語}という．
そして，述語を用いて命題を表現する論理を\emph{述語論理}\index{じゅつごろんり@述語論理}という．


述語論理においてはいくつかの記号を用いる．次の表で確認しよう：

\begin{table}[ht]
  \centering
  \caption{述語論理の記号}
  \begin{tabular}{c|c}
    \hline
    記号        & 意味                              \\
    \hline
    $\forall$ & 任意の（全称量化\index{ぜんしょうりょうか@全称量化}） \\
    $\exists$ & 存在する（存在量化\index{そんざいりょうか@存在量化}） \\
    \hline
  \end{tabular}
\end{table}

これらの記号を用いて，述語を命題にすることができる．$\forall$と$\exists$をまとめて\emph{量化子}\index{りょうかし@量化子}とよぶ．

\phantomsection
\begin{definition}{全称命題}{全称命題}
  $P(x)$を述語とする．
  \[
    \forall\, x \, : \, P(x)
  \]
  は，「すべての$x$について，$P(x)$が成り立つ」という意味であり，これを\emph{全称命題}\index{ぜんしょうめいだい@全称命題}という．
\end{definition}

\phantomsection
\begin{definition}{存在命題}{存在命題}
  $Q(x)$を述語とする．
  \[
    \exists \,  x\, : \, Q(x)
  \]
  は，「$Q(x)$が成り立つような$x$が存在する」という意味であり，これを\emph{存在命題}\index{そんざいめいだい@存在命題}という．
\end{definition}

\phantomsection
\begin{example}{量化記号の例}{量化記号の例}
  \begin{itemize}
    \item $\forall x \in \mathbb{R}\, : \, x^2 \geqq 0$\\
          「任意の実数$x$について，$x^2$は$0$以上である．」
    \item $\exists x \in \mathbb{Z} \, : \, x^2 = 4$\\
          「ある整数$x$が存在して，$x^2 = 4$となる．」
  \end{itemize}
\end{example}


\subsection{入れ子の量化子\index{いれこのりょうかし@入れ子の量化子}}
 
ここまでは変数がひとつの述語を扱ってきたが，$P(x, y)$のように変数を2つ以上含む述語も考えられる．
このとき，命題を作るにはすべての変数を量化する必要があり，量化子が
\[
  \forall x,\, \exists y \, : \, P(x, y)
\]
のように2つ以上並ぶことになる．これを量化子の\emph{入れ子}\index{いれこのりょうかし@入れ子の量化子}という．
 
入れ子になった量化子は，\textbf{左から順に}読む．
すなわち，上の例は「任意の$x$に対して，ある$y$が存在して，$P(x, y)$が成り立つ」と読む．
このとき重要なのは，\textbf{内側（右側）の変数は，外側（左側）の変数に依存してよい}ということである．
上の例でいえば，$y$は$x$に応じて選び直してよい．
 
\phantomsection
\subsubsection{同種の量化子の入れ子}
 
まず，同じ種類の量化子が並ぶ場合を考えよう．この場合，量化子の順序を入れ替えても命題の意味は変わらない：
 
\begin{prop}{同種の量化子の交換}{同種の量化子の交換}
  $P(x, y)$を2変数の述語とする．このとき，次が成り立つ：
  \begin{align*}
    \forall x,\, \forall y \, : \, P(x, y) & \equiv \forall y,\, \forall x \, : \, P(x, y) \\
    \exists x,\, \exists y \, : \, P(x, y) & \equiv \exists y,\, \exists x \, : \, P(x, y)
  \end{align*}
\end{prop}
 
直観的には，$\forall x,\, \forall y$は「どの$x$とどの$y$の組合せについても」，
$\exists x,\, \exists y$は「うまい$x$と$y$の組がひとつでもあれば」という意味であり，
いずれも$x$と$y$のどちらを先に述べるかは関係ない．
 
\phantomsection
\begin{example}{同種の量化子の例}{同種の量化子の例}
  \begin{itemize}
    \item $\forall x \in \mathbb{R},\, \forall y \in \mathbb{R} \, : \, x^2 + y^2 \geqq 2xy$\\
          「任意の実数$x$と任意の実数$y$について，$x^2 + y^2 \geqq 2xy$が成り立つ．」\\
          これは$x^2 + y^2 - 2xy = (x - y)^2 \geqq 0$からわかる．
          $x$と$y$の役割を入れ替えても主張は変わらない．
    \item $\exists x \in \mathbb{Z},\, \exists y \in \mathbb{Z} \, : \, x^2 + y^2 = 5$\\
          「ある整数$x$とある整数$y$が存在して，$x^2 + y^2 = 5$となる．」\\
          実際，$(x, y) = (1, 2)$がその一例である．
          こちらも，$x$と$y$のどちらを先に見つけるかは問題にならない．
  \end{itemize}
\end{example}
 
\phantomsection
\subsubsection{異種の量化子の入れ子}
 
一方，$\forall$と$\exists$が混在する場合には，順序が本質的に重要になる．
\[
  \forall x,\, \exists y \, : \, P(x, y)
\]
は「任意の$x$に対して，（$x$に応じた）$y$が存在する」という意味であり，$y$は$x$ごとに選び直してよい．
これに対して，
\[
  \exists y,\, \forall x \, : \, P(x, y)
\]
は「ある$y$が存在して，その$y$は任意の$x$に対して$P(x, y)$を満たす」という意味であり，
$y$は$x$によらず\textbf{ひとつに固定}されていなければならない．
後者のほうが強い主張であり，一般に次の一方向の含意のみが成り立つ：
 
\begin{prop}{異種の量化子の順序}{異種の量化子の順序}
  $P(x, y)$を2変数の述語とする．このとき，次が成り立つ：
  \[
    \bigl( \exists y,\, \forall x \, : \, P(x, y) \bigr)
    \to
    \bigl( \forall x,\, \exists y \, : \, P(x, y) \bigr)
  \]
  この逆は，一般には成り立たない．
\end{prop}
 
すなわち，「全員に共通のうまい$y$がある」ならば「各人に応じたうまい$y$がある」のは当然だが，
その逆は保証されない．逆が成り立たない例として，次のアルキメデスの原理を見てみよう．
 
学部1年の微分積分学の講義で，アルキメデスの原理
\[
  \forall a>0,\, \exists n\in\mathbb{N} : n>a
\]
を学ぶことがある．この命題を学習する際には，$\forall$と$\exists$が入れ子になっていることに注意しよう．
アルキメデスの原理の主張の解釈は「任意の正の実数$a$に対して，それより大きい自然数$n$が存在する」
である．これを
\[
  \exists n\in\mathbb{N},\, \forall a>0 : n>a
\]
とかき換えると，意味は大きく変わってしまう．
こちらは「ある自然数$n$が存在して，すべての正の実数$a$より大きい」
という意味になるが，そのような自然数は存在しないので偽である．
 
このように，量化記号の順序を入れ替えると，
命題の意味だけでなく真偽まで変わることがある．
$\varepsilon$-$\delta$論法では，
\[
  \forall \varepsilon>0,\, \exists \delta>0
\]
という形で量化記号が入れ子になっており，
$\delta$は選んだ$\varepsilon$に応じて決めてよいことを表している．
 
\phantomsection
\subsubsection{入れ子の量化子の否定}
 
最後に，量化子を含む命題の否定について述べる．
量化子の否定は，$\lnot$が量化子を通過するたびに$\forall$と$\exists$が入れ替わる，という規則に従う：
 
\begin{prop}{量化子の否定}{量化子の否定}
  $P(x)$を述語とする．このとき，次が成り立つ：
  \begin{align*}
    \lnot \bigl( \forall x \, : \, P(x) \bigr) & \equiv \exists x \, : \, \lnot P(x) \\
    \lnot \bigl( \exists x \, : \, P(x) \bigr) & \equiv \forall x \, : \, \lnot P(x)
  \end{align*}
\end{prop}
 
これはド・モルガンの法則の量化子版とみなせる．
実際，「すべての$x$で$P(x)$が成り立つ」の否定は「$P(x)$が成り立たない$x$が（少なくともひとつ）存在する」であり，
「$P(x)$が成り立つ$x$が存在する」の否定は「すべての$x$で$P(x)$が成り立たない」である．
 
この規則を入れ子の量化子に対して\textbf{外側から順に}適用すれば，入れ子になった命題の否定も機械的に作ることができる：
\begin{align*}
  \lnot \bigl( \forall x,\, \exists y \, : \, P(x, y) \bigr)
  & \equiv \exists x,\, \lnot \bigl( \exists y \, : \, P(x, y) \bigr) \\
  & \equiv \exists x,\, \forall y \, : \, \lnot P(x, y)
\end{align*}
すなわち，量化子の列$\forall x,\, \exists y$を$\exists x,\, \forall y$にすべて反転させ，
最後に内側の述語を否定すればよい．
この操作は，$\varepsilon$-$\delta$論法で「収束しないこと」を定式化する際などに繰り返し現れる．


% NOTE: この部は執筆中（今後，節を追記する可能性あり）
